\documentclass{habshe-tufte}

\usepackage{lipsum}
\usepackage{amsmath,amssymb,amsfonts}

\begin{document} % implied full-width layout

\subject{Subject Huge}

\heading{Heading huge}

\subheading{Sub-Heading LARGE}

\subsubheading{Sub-Sub-Heading Large}

\subsubsubheading{Sub-Sub-Sub-Heading large}

Regular paragraphs are in normal size, indented with no vertical space between them.

\lipsum[4-6]

\begin{topic}[Topic] % run-in: first sentence follows on the SAME line (no blank line)
  Topic environments start with an unindented line where the topic name starts the line and the first paragraph continues after it inline. Subsequent paragraphs indent with no vertical space.

  \lipsum[7-8]

  \begin{subtopic}[Sub-Topic]
    Subtopic environments are a regular indented paragraph same as a regular paragraph in topic, but start with a subtopic name. The first paragraph continues after the subtopic name inline. Subsequent paragraphs are indented.

    \lipsum[9-10]
  \end{subtopic}

  \lipsum[11]
\end{topic}

Sometimes math is inline:

\begin{itemize}
\item The scalar value $\alpha \in \mathbb{R}$ satisfies the condition $\alpha^2 > 0$ for all non-zero inputs.
\item Let $f(x) = e^{-x} \sin(x)$ be a continuous mapping on the open interval $(0, 1)$.
\item Given vectors $\mathbf{u} = (1, 2, 3)$ and $\mathbf{v} = (4, 5, 6)$, their dot product is $\mathbf{u} \cdot \mathbf{v} = 32$.
\item The function $g(z) = \frac{1}{2\pi i} \oint_{\gamma} \frac{g(w)}{w - z} \, dw$ defines the Cauchy integral formula.
\item Lorem ipsum dolor sit amet, ubi $\sum_{n=1}^{\infty} a_n$ converges absolutely.
\end{itemize}

Other times it is not:

$$
ax^2 + bx = -c
$$

Divide both sides by $a$:

$$
x^2 + \frac{b}{a}x = -\frac{c}{a}
$$

Add $\left(\frac{b}{2a}\right)^2$ to both sides to complete the square:

$$
x^2 + \frac{b}{a}x + \left(\frac{b}{2a}\right)^2 = \frac{b^2}{4a^2} - \frac{c}{a}
$$

$$
\left(x + \frac{b}{2a}\right)^2 = \frac{b^2 - 4ac}{4a^2}
$$

Taking the square root of both sides gives:

$$
x + \frac{b}{2a} = \pm \frac{\sqrt{b^2 - 4ac}}{2a}
$$

Solving for $x$, we arrive at the quadratic formula:

$$
x = \frac{-b \pm \sqrt{b^2 - 4ac}}{2a}
$$

Proof: Sum of Consecutive Integers
\textbf{Theorem:} For any positive integer $n$, the sum of the first $n$ natural integers is given by:

$$
\sum_{i=1}^{n} i = \frac{n(n+1)}{2}
$$

\begin{enumerate}
\item \textbf{Base Case:} For $n = 1$, the left-hand side is $1$, and the right-hand side is:
  $$
  \frac{1(1+1)}{2} = \frac{2}{2} = 1
  $$
  Thus, the base case holds.

\item \textbf{Inductive Hypothesis:} Assume the formula holds for some integer $k \ge 1$:
  $$
  \sum_{i=1}^{k} i = \frac{k(k+1)}{2}
  $$

\item \textbf{Inductive Step:} We must show that it holds for $n = k + 1$:
  $$
  \sum_{i=1}^{k+1} i = \left( \sum_{i=1}^{k} i \right) + (k+1)
  $$

  Substitute the inductive hypothesis into the equation:
  $$
  \sum_{i=1}^{k+1} i = \frac{k(k+1)}{2} + (k+1)
  $$

  Factor out $(k+1)$:
  $$
  \sum_{i=1}^{k+1} i = (k+1) \left( \frac{k}{2} + 1 \right)
  $$

  Combine terms inside the parenthesis:
  $$
  \sum_{i=1}^{k+1} i = (k+1) \left( \frac{k + 2}{2} \right) = \frac{(k+1)(k+2)}{2}
  $$

  This matches the formula for $n = k + 1$. Therefore, by mathematical induction, the statement is true for all $n \in \mathbb{N}$. $\blacksquare$
\end{enumerate}

\tufte

\subject{Subject Huge}

\heading{Heading huge}

\subheading{Sub-Heading LARGE}

\subsubheading{Sub-Sub-Heading Large}

\subsubsubheading{Sub-Sub-Sub-Heading large}

Regular paragraphs are in normal size, indented with no vertical space between them.

\lipsum[4-6]

\begin{topic}[Topic]
  Topic environments start with an unindented line where the topic name starts the line and the first paragraph continues after it inline. Subsequent paragraphs indent with no vertical space.

  \lipsum[7-8]

  \begin{subtopic}[Sub-Topic]
    Subtopic environments are a regular indented paragraph same as a regular paragraph in topic, but start with a subtopic name. The first paragraph continues after the subtopic name inline. Subsequent paragraphs are indented.

    \lipsum[9-10]
  \end{subtopic}

  \lipsum[11]
\end{topic}

Sometimes math is inline:

\begin{itemize}
\item The scalar value $\alpha \in \mathbb{R}$ satisfies the condition $\alpha^2 > 0$ for all non-zero inputs.
\item Let $f(x) = e^{-x} \sin(x)$ be a continuous mapping on the open interval $(0, 1)$.
\item Given vectors $\mathbf{u} = (1, 2, 3)$ and $\mathbf{v} = (4, 5, 6)$, their dot product is $\mathbf{u} \cdot \mathbf{v} = 32$.
\item The function $g(z) = \frac{1}{2\pi i} \oint_{\gamma} \frac{g(w)}{w - z} \, dw$ defines the Cauchy integral formula.
\item Lorem ipsum dolor sit amet, ubi $\sum_{n=1}^{\infty} a_n$ converges absolutely.
\end{itemize}

Other times it is not:

$$
ax^2 + bx = -c
$$

Divide both sides by $a$:

$$
x^2 + \frac{b}{a}x = -\frac{c}{a}
$$

Add $\left(\frac{b}{2a}\right)^2$ to both sides to complete the square:

$$
x^2 + \frac{b}{a}x + \left(\frac{b}{2a}\right)^2 = \frac{b^2}{4a^2} - \frac{c}{a}
$$

$$
\left(x + \frac{b}{2a}\right)^2 = \frac{b^2 - 4ac}{4a^2}
$$

Taking the square root of both sides gives:

$$
x + \frac{b}{2a} = \pm \frac{\sqrt{b^2 - 4ac}}{2a}
$$

Solving for $x$, we arrive at the quadratic formula:

$$
x = \frac{-b \pm \sqrt{b^2 - 4ac}}{2a}
$$

Proof: Sum of Consecutive Integers
\textbf{Theorem:} For any positive integer $n$, the sum of the first $n$ natural integers is given by:

$$
\sum_{i=1}^{n} i = \frac{n(n+1)}{2}
$$

\begin{enumerate}
\item \textbf{Base Case:} For $n = 1$, the left-hand side is $1$, and the right-hand side is:
  $$
  \frac{1(1+1)}{2} = \frac{2}{2} = 1
  $$
  Thus, the base case holds.

\item \textbf{Inductive Hypothesis:} Assume the formula holds for some integer $k \ge 1$:
  $$
  \sum_{i=1}^{k} i = \frac{k(k+1)}{2}
  $$

\item \textbf{Inductive Step:} We must show that it holds for $n = k + 1$:
  $$
  \sum_{i=1}^{k+1} i = \left( \sum_{i=1}^{k} i \right) + (k+1)
  $$

  Substitute the inductive hypothesis into the equation:
  $$
  \sum_{i=1}^{k+1} i = \frac{k(k+1)}{2} + (k+1)
  $$

  Factor out $(k+1)$:
  $$
  \sum_{i=1}^{k+1} i = (k+1) \left( \frac{k}{2} + 1 \right)
  $$

  Combine terms inside the parenthesis:
  $$
  \sum_{i=1}^{k+1} i = (k+1) \left( \frac{k + 2}{2} \right) = \frac{(k+1)(k+2)}{2}
  $$

  This matches the formula for $n = k + 1$. Therefore, by mathematical induction, the statement is true for all $n \in \mathbb{N}$. $\blacksquare$
\end{enumerate}

\end{document}
